% ============================================================
%  APPENDIX TO CHAPTER 1 -- EXERCISES S1.1-S1.15
%  Imported from Tao I, ECNU, Yu Pin.  Not Rudin's.
% ============================================================
\ssection{Appendix exercises}

\begin{roadmap}[How these differ from Rudin's twenty]
Rudin's own Exercises 1.1--1.20 end the chapter proper. They are almost all
verifications inside a structure already given: check a field axiom, compute
with cuts, establish an inequality. These fifteen come from the comparison
texts and ask two other kinds of question --- what the axioms \emph{cost},
and how to actually compute with them.

\medskip
\raggedright
\textbf{Group I (1--5): the toolkit.} Drills on the four facts of Supplement
D, plus the two supremum lemmas every later chapter uses and Rudin never
assigns. Exercise 4 is the first place in this book where anyone is asked to
\emph{compute} a supremum.

\medskip
\textbf{Group II (6--10): what the axioms buy.} Where each hypothesis is
spent --- induction failing on $\Z$, the supremum principle producing
Archimedes, and the disconcerting question of whether one can prove any
real numbers exist at all.

\medskip
\textbf{Group III (11--15): decimals and construction.} Exactly which reals
have two decimal names, the supremum built by hand, and a game.

\medskip
\textbf{Marked $\star$:} harder, or requiring a result from a later chapter.
\end{roadmap}

\begin{enumerate}[label=\textbf{S1.\arabic*.}, leftmargin=3.4em, itemsep=9pt,
                  topsep=8pt, labelsep=0.5em]

% ---------------- Group I ----------------
\item Let $a, b \in \R$. Prove that if $|a-b| < \varepsilon$ for every
$\varepsilon > 0$, then $a = b$.\kn{The $\varepsilon$-principle of S1.9 as a
drill, and the single most-used move in the book. The proof is one line by
contradiction --- take $\varepsilon = |a-b|$ --- but the habit it builds is
what matters: to prove two things equal, show their distance is smaller than
every positive number. Rudin uses this at 1.11, 1.21, 3.2(b) and throughout
Chapter 3 without ever stating it.}
\provenance{ECNU \cjk{习题}{Ex} 1.1 \#3 (bk.\ p.\,4).}

\item Prove that for all $x \in \R$: (a) $|x-1| + |x-2| \ge 1$; (b)
$|x-1| + |x-2| + |x-3| \ge 2$. Determine in each case exactly when equality
holds.\kn{Read each expression as a total distance to fixed points on the
line. In (a) the sum of distances to $1$ and $2$ is least --- and equal to
the gap --- exactly when $x$ lies between them. In (b) the middle point is
free to be hit exactly, so equality occurs only at $x=2$: the minimiser of a
sum of distances is the median. A geometric habit worth acquiring before
Chapter 2 makes distance abstract.}
\provenance{ECNU \cjk{习题}{Ex} 1.1 \#5 (bk.\ p.\,4).}

\item Let $a,b,c > 0$. Prove
$\bigl|\sqrt{a^2+b^2} - \sqrt{a^2+c^2}\bigr| \le |b-c|$, and give the
inequality a geometric meaning.\kn{Two right triangles sharing the leg $a$:
the hypotenuses cannot differ by more than the other legs do. Algebraically
it says $t \mapsto \sqrt{a^2+t^2}$ never stretches distances --- it is what
Chapter 4 will call a Lipschitz function with constant $1$. Rudin's Chapter 1
exercises contain nothing geometric; this one rewards a picture.}
\provenance{ECNU \cjk{习题}{Ex} 1.1 \#6 (bk.\ p.\,4).}

\item For each set, find $\sup$ and $\inf$ and verify the answer
\emph{from Definition 1.8}: (a) $\{x : x^2 < 2\}$; (b) $\{n! : n \in
\N\}$; (c) $\{x \text{ irrational} : 0<x<1\}$; (d)
$\{1 - 2^{-n} : n \in \N\}$.\kn{Rudin's Chapter 1 never asks anyone to
compute a supremum, and computing four of them is the fastest way to
internalise that Definition 1.8 has \emph{two} clauses --- an upper bound,
and no smaller one. Note (b), where the set is unbounded above and the
answer is $+\infty$ only in the extended sense of 1.23; and (c), where the
supremum is $1$ and is not attained, while every neighbourhood of it is
crowded.}
\provenance{ECNU \cjk{习题}{Ex} 1.2 \#4 (bk.\ p.\,8).}

\item Let $A,B$ be nonempty and bounded, and $A+B = \{x+y : x \in A,\ y \in
B\}$. Prove $\sup(A+B) = \sup A + \sup B$ and
$\inf(A+B) = \inf A + \inf B$.\kn{The workhorse lemma behind every
limit-algebra proof, and it is not in Rudin's exercise set. Both halves are
needed: ``$\le$'' because $\sup A + \sup B$ bounds $A+B$, and ``$\ge$''
because any smaller bound would fail for some $x$ near $\sup A$ and $y$ near
$\sup B$ --- which is the second clause of 1.8 doing its work. Compare 3.3,
where the same statement for limits gets a whole theorem.}
\provenance{ECNU \cjk{习题}{Ex} 1.2 \#7 (bk.\ p.\,8).}

% ---------------- Group II ----------------
\item Show that the principle of induction does not apply to $\Z$: exhibit a
property $P(n)$ of integers such that $P(0)$ holds and $P(n)$ implies
$P(n+1)$ for every integer $n$, yet $P(n)$ fails for some
integer.\kn{Take $P(n)$ to be ``$n \ge 0$.'' One line, and a permanent
lesson: induction is not a general logical principle but a specific axiom
about $\N$, and what makes it work is that $\N$ has a starting point with
nothing before it. Rudin uses induction freely throughout Chapter 1 --- at
1.1, at 1.21, in the appendix --- and never says what licenses it.}
\provenance{Tao I, Ex.\ 4.1.8 (bk.\ p.\,81).}

\item Assuming only the field axioms and the order axioms, prove that the
least-upper-bound property implies the archimedean property.\kn{This is
Rudin's Theorem 1.20(a), but posed inside an axiom system where Archimedes is
\emph{independent}, so that the exercise shows what is being bought rather
than merely how. Compare S1.3, and note the shape of the argument: if $\N$
were bounded, $\sup\N - 1$ would fail to be an upper bound, and the integer
beating it has a successor beating the supremum.}
\provenance{Yu Pin, homework A7 (p.\,36).}

\item $\star$~Suppose $a>0$. Can you prove that the open interval $(0,a)$ is
nonempty?\kn{Disconcerting, and meant to be. On the field and order axioms
alone the only elements you have names for are $0$, $1$ and $-1$, so
producing a real strictly between $0$ and $a$ is not the triviality it looks
like --- you must manufacture new elements, and manufacturing elements in
bulk is precisely what the completeness axiom is for. Rudin's 1.19 hands the
reader a fully populated $\R$ by fiat, which is why the question cannot even
be asked inside his development. Yu Pin's own comment: ``you will see the
shadow of the Dedekind cut here.''}
\provenance{Yu Pin, \cjk{思考题}{thought question} (p.\,22).}

\item Prove that the ordered field $\R(t)$ of Supplement B really is ordered
--- that the sum and product of eventually-positive rational functions are
eventually positive, and that trichotomy holds --- and then verify directly
that the set $\N$ has no least upper bound in it.\kn{The verification behind
S1.2 and S1.3. For trichotomy, a nonzero rational function has finitely many
zeros, so it has constant sign near $+\infty$. For the second part, show that
if $f$ bounded $\N$ above then so would $f-1$, since $f - 1 > n$ whenever
$f > n+1$ --- so no candidate can be least. The example is the one thing
Chapter 1 lacks: a field where Archimedes visibly fails.}
\provenance{Supplied here; see the provenance note at S1.2.}

\item $\star$~Prove that, given the field and order axioms together with the
archimedean property, the least-upper-bound property and the nested-interval
property are equivalent. Then explain why the archimedean hypothesis cannot
be dropped from this statement.\kn{The precise form of Chapter 1's
comparison box. One direction is a bisection; the other is the argument of
S2.3 in the Chapter 2 appendix. For the last part, look at S1.2: nested
intervals alone cannot rule out infinite elements, so it cannot imply
Archimedes, whereas \lub{} implies both at once. This is exactly why Yu Pin
lists four axioms where Rudin needs one.}
\provenance{Yu Pin, Remark 4 (p.\,28) and Thm 5 (pp.\,26--27).}

% ---------------- Group III ----------------
\item Prove that a real number has two decimal representations if it is a
terminating decimal --- that is, if $x = n/10^m$ for integers $n,m$ --- and
exactly one otherwise. Show also that the only decimal representations of
$1$ are $1.000\ldots$ and $0.999\ldots$\kn{The statement Rudin's 1.22
declines to make. He asserts that the expansion exists and is unique, having
silently fixed the non-terminating digits by always choosing the
\emph{largest} digit at each stage; ECNU fixes the opposite convention
(S1.4). The exercise is to see that the ambiguity is real, is confined to
exactly the terminating decimals, and is always a choice between a
$0$-tail and a $9$-tail.}
\provenance{Tao I, Ex.\ B.2.2--B.2.3 (bk.\ p.\,338); Yu Pin, homework E1
(p.\,506), which states it as a four-way equivalence.}

\item Prove the supremum principle by the digit construction of S1.7 for a
set of your own choosing --- for instance $S = \{x : x^2 < 2\}$ --- and
carry the algorithm far enough to produce the first four digits of
$\sup S$.\kn{The point is to watch Definition 1.8's two clauses act as loop
invariants: at each stage one keeps the largest digit that some element still
reaches. Doing it once on $\{x : x^2<2\}$ produces $1.414$ and makes
concrete what Rudin's 1.1 says is missing from $\Q$ --- the digits are
perfectly well defined at every finite stage, and only the infinite string
fails to be rational.}
\provenance{ECNU \cjk{定理}{Thm} 1.1 (bk.\ pp.\,5--7).}

\item $\star$~Construct the supremum of a nonempty bounded set by successive
rational approximations, \emph{without} invoking the least-upper-bound
property: for each $n$ find the unique integer $m$ such that $m/n$ is an
upper bound while $(m-1)/n$ is not, and prove $m$ is unique.\kn{The same
construction used in S2.7 to recover $\sup$ from Cauchy completeness, here
isolated. Uniqueness of $m$ is where the archimedean property enters --- it
is what guarantees the set of integers $k$ with $k/n$ an upper bound is
nonempty and bounded below. Compare with exercise 12: one construction walks
in steps of $10^{-k}$, the other in steps of $1/n$.}
\provenance{Tao I, Ex.\ 5.5.2--5.5.3 (bk.\ p.\,121).}

\item Prove that $b^{m/n}$ is well defined --- that
$(b^{p})^{1/q} = (b^{m})^{1/n}$ whenever $m/n = p/q$ --- and deduce the law
$b^{r+s} = b^rb^s$ for rational $r,s$.\kn{Rudin's Exercise 1.6, and the
independence step is the whole of it. Both sides are positive $n q$th roots
of $b^{mq}$, so the \emph{uniqueness} half of 1.21 identifies them --- one
more instance of the pattern flagged at 1.21, where an existence-and-
uniqueness theorem is used entirely for its uniqueness. See S1.13.}
\provenance{Tao I, Lemma 5.6.8 (bk.\ p.\,124); Rudin's Ex.\ 1.6.}

\item $\star$~\emph{(A game.)} Two players alternately choose nested closed
intervals $I_1 \supset I_2 \supset \cdots$, each at most half the length of
its predecessor, so that the intersection is a single point $x$. The first
player wins if $x$ is rational, the second if $x$ is irrational. Who has a
winning strategy?\kn{Everything needed lives in Chapter 1: nested intervals,
density of $\Q$, and the uncountability that follows from S2.17 in the
Chapter 2 appendix. The second player wins, and the reason is the same
reason the diagonal argument works --- the first player must commit to a
countable list of targets while the second retains room to dodge each one in
turn. A pleasant way to see that ``most'' reals are irrational without ever
using the word measure.}
\provenance{Yu Pin, problem G (p.\,69), the Banach--Mazur game.}

\end{enumerate}
